% =====================================================================
%  Apéndices
% =====================================================================
\chapter{Mensajes del sistema}\label{ap:mensajes}

\entradilla{Referencia rápida de los avisos que muestra el panel después de cada acción,
agrupados por módulo.}

\begin{longtblr}{colspec={Q[3.1cm,l] X[l] Q[4.2cm,l]}}
  Módulo & Mensaje & Cuándo aparece \\
  Cotizaciones & \msg{URL de cotización subida exitosamente.} & Al enviar la propuesta. \\
  Cotizaciones & \msg{La URL de la cotización es obligatoria.} & Al enviar sin enlace. \\
  Cotizaciones & \msg{La URL debe ser una dirección válida.} & El enlace no tiene formato de dirección web. \\
  Clientes & \msg{No se pudo iniciar la exportación. Intente nuevamente.} & Falla al iniciar el Excel. \\
  Clientes & \msg{No se pudo generar el archivo. Intente nuevamente.} & Falla al generar el Excel. \\
  Clientes & \msg{Se perdió la conexión durante la exportación.} & Se cortó internet mientras se generaba. \\
  Clientes & \msg{La generación está tardando demasiado. Verifique que el worker de la cola esté en ejecución.} & El proceso en segundo plano no responde. \\
  Unidades de negocio & \msg{La unidad de negocio «\var{nombre}» ha sido creada exitosamente.} & Al crear. \\
  Unidades de negocio & \msg{La unidad de negocio «\var{nombre}» ha sido actualizada exitosamente.} & Al editar. \\
  Unidades de negocio & \msg{La unidad de negocio «\var{nombre}» ha sido eliminada exitosamente.} & Al eliminar. \\
  Unidades de negocio & \msg{No se puede eliminar la unidad de negocio «\var{nombre}» porque tiene servicios asociados.} & Al eliminar una unidad con servicios. \\
  Líneas de gestión & \msg{Línea de gestión creada correctamente.} & Al crear. \\
  Líneas de gestión & \msg{Línea de gestión actualizada exitosamente.} & Al editar. \\
  Líneas de gestión & \msg{Línea de gestión eliminada exitosamente.} & Al eliminar. \\
  Líneas de gestión & \msg{No se puede eliminar la línea de gestión «\var{nombre}» porque tiene servicios asociados.} & Al eliminar una línea con servicios. \\
  Servicios & \msg{Servicio creado correctamente.} & Al crear. \\
  Servicios & \msg{Servicio actualizado exitosamente.} & Al editar. \\
  Servicios & \msg{Servicio eliminado exitosamente.} & Al eliminar. \\
  Servicios & \msg{No existen unidades de negocio con las cuales usted pueda asociar un servicio.} & Al crear un servicio sin unidades. \\
  Servicios & \msg{No existen líneas de gestión con las cuales usted pueda asociar un servicio.} & Al crear un servicio sin líneas. \\
  Árbol de decisión & \msg{Pregunta «\var{pregunta}» creada.} & Al crear una pregunta. \\
  Árbol de decisión & \msg{Pregunta «\var{pregunta}» guardada.} & Al editar una pregunta. \\
  Árbol de decisión & \msg{Pregunta «\var{pregunta}» eliminada.} & Al eliminar una pregunta. \\
  Árbol de decisión & \msg{Preguntas copiadas desde «\var{unidad}».} & Al copiar un cuestionario. \\
  Árbol de decisión & \msg{Tiene cambios sin guardar en la pregunta que está editando. ¿Desea descartarlos?} & Al salir de una pregunta con cambios. \\
  Usuarios & \msg{Usuario creado correctamente.} & Al crear. \\
  Usuarios & \msg{Usuario actualizado correctamente.} & Al editar. \\
  Usuarios & \msg{Usuario eliminado correctamente.} & Al eliminar. \\
  Usuarios & \msg{No puedes eliminar tu propio usuario.} & Al intentar borrar la cuenta propia. \\
  Correos de bienvenida & \msg{Se pusieron en cola \var{N} correos de bienvenida.} & Al enviar la lista. \\
  Correos de bienvenida & \msg{Se omitieron \var{X} líneas sin correo válido (ej.: …).} & Si la lista tenía líneas inválidas. \\
  Correos de bienvenida & \msg{Máximo 500 correos por envío. Reduzca la lista e intente nuevamente.} & Lista de más de 500 correos. \\
  Configuración & \msg{Configuración actualizada exitosamente.} & Al guardar. \\
\end{longtblr}

% ---------------------------------------------------------------------
\chapter{Preguntas frecuentes}\label{ap:faq}

\entradilla{Respuestas a las situaciones más comunes y qué hacer en cada caso.}

\newcommand{\faq}[2]{%
  \begin{tcolorbox}[enhanced, colback=white, colframe=linea, boxrule=0.8pt, arc=6pt,
      left=12pt, right=12pt, top=9pt, bottom=9pt, before skip=8pt, after skip=8pt]
    {\montserrat\bfseries\color{azul}{\iconos\char"F059}\hspace{0.5em}#1}\par\vspace{3pt}
    #2
  \end{tcolorbox}}

\faq{Creé una unidad de negocio, pero no aparece en Training~InnovaQuote.}
  {Una unidad solo se muestra a los clientes cuando su árbol de decisión tiene al menos una
   pregunta. Cree sus preguntas o copie las de otra unidad (\seccion{sec:copiar-arbol}).}

\faq{No puedo eliminar una unidad de negocio o una línea de gestión.}
  {Tiene servicios asociados. Elimine esos servicios o asígnelos a otra unidad o línea
   (\seccion{sec:servicios}) e intente de nuevo.}

\faq{Envié la propuesta con un enlace equivocado.}
  {El panel no permite cambiar el enlace después de enviarlo. Envíe al cliente el enlace
   correcto por correo y pida al equipo técnico que corrija el registro, indicando el código de
   la cotización.}

\faq{El cliente dice que no recibió el correo con la propuesta.}
  {Pídale que revise la carpeta de correo no deseado. Verifique en la ficha del cliente que su
   correo esté bien escrito (\seccion{sec:clientes}). Los correos se envían en segundo plano y
   pueden tardar unos minutos; si no llegan, avise al equipo técnico.}

\faq{No recibo los avisos de nuevas solicitudes.}
  {Revise que el \campo{Correo de notificaciones} esté diligenciado en
   \menu{Sistema}{Configuración} (\seccion{sec:configuracion}) y revise la carpeta de correo no
   deseado de ese buzón.}

\faq{El cliente dice que adjuntó documentos, pero no los encuentro en el panel.}
  {El panel no guarda los archivos: solo muestra sus nombres en la respuesta \campo{Adjuntar
   documentos}. Los archivos viajan adjuntos en el correo \msg{Nueva solicitud de cotización
   \var{código}} que llega al \campo{Correo de notificaciones}. Búsquelo en ese buzón
   (\seccion{sec:documentos-cliente}). Si ese correo estaba vacío cuando llegó la solicitud,
   los documentos se descartaron y hay que pedírselos de nuevo al cliente.}

\faq{La exportación a Excel se queda en «Generando…».}
  {El archivo se genera con un proceso en segundo plano del servidor. Si después de unos
   minutos aparece \msg{La generación está tardando demasiado…}, avise al equipo técnico para
   que revise ese proceso.}

\faq{El botón «Crear pregunta» o «Guardar pregunta» está desactivado.}
  {Hay un problema en la pregunta o no ha hecho cambios. Revise el recuadro rojo dentro de la
   tarjeta: indica lo que debe corregir (\seccion{sec:validaciones-arbol}).}

\faq{Una pregunta de rama tiene el borde ámbar.}
  {Ninguna opción lleva a ella, así que el cliente no la verá. Conéctela desde una opción con
   \campo{Al elegirla} o elimínela (\seccion{sec:ramas}).}

\faq{Olvidé mi contraseña.}
  {Use \textbf{¿Olvidó su contraseña?} en la pantalla de inicio de sesión
   (\seccion{sec:recuperar}) o pida a otro administrador que le asigne una nueva
   (\seccion{sec:usuario-editar}).}

\faq{Necesito retirar el acceso a una persona.}
  {Elimine su cuenta en \menu{Sistema}{Usuarios}. Si la cuenta es la suya, pídaselo a otro
   administrador.}

% ---------------------------------------------------------------------
\chapter{Glosario}\label{ap:glosario}

\entradilla{Términos que se usan en el panel y en este manual.}

\begin{tabla}{colspec={Q[4.2cm,l] X[l]}}
  Término & Definición \\
  Abreviado & Letras y números que identifican a una unidad de negocio y forman el inicio de
    sus códigos de cotización, por ejemplo \textbf{EO}. \\
  Árbol de decisión & El conjunto de preguntas que responde el cliente de una unidad al cotizar,
    con sus opciones y ramas. \\
    Bloque del formulario & Cada uno de los cuatro grupos temáticos en que se organizan las
    preguntas del formulario. \\
  Cliente & Persona que ha enviado al menos una solicitud desde Training~InnovaQuote. Se identifica por
    su correo. \\
  Código de cotización & Identificador único de cada solicitud, con la forma
    \textbf{ABREVIADO-V-AAnn}, por ejemplo \textbf{EO-V-2601}. \\
  Correo de notificaciones & Buzón interno que recibe un aviso con cada solicitud nueva y cada
    propuesta enviada. \\
    Formulario de solicitud & El formulario que diligencia el cliente al solicitar una cotización.
    Se configura en el árbol de decisión. \\
  Línea de gestión & Agrupación temática de los servicios, por ejemplo HSEQ. \\
  Orden principal & La secuencia numerada de preguntas que todos los clientes responden. \\
    Panel de cotizaciones & La parte privada de Training~InnovaQuote, donde trabaja el administrador. \\
  
  Pregunta de rama & Pregunta que solo aparece cuando el cliente elige cierta opción. \\
  Propuesta & El documento con la oferta comercial que se envía al cliente mediante un enlace. \\
  
  Servicio & Un servicio de la oferta de una unidad de negocio, clasificado en una línea de
    gestión. \\
    Training~InnovaQuote & Nombre de la plataforma web de cotizaciones de Training Corporation, disponible en
    \ruta{\dominio}. \\
  Unidad de negocio & Tipo de servicio que el cliente elige al cotizar, por ejemplo Auditoría,
    Consultoría o Formación. \\
  Usuarios simultáneos & Personas que usan Training~InnovaQuote al mismo tiempo. \\
\end{tabla}
